\documentclass[11pt]{article}
\usepackage[margin=1in]{geometry}
\usepackage{amsmath,amssymb,amsthm,bm,mathtools}
\usepackage{hyperref}
\newtheorem{theorem}{Theorem}
\newtheorem{lemma}{Lemma}
\newtheorem{corollary}{Corollary}
\theoremstyle{definition}
\newtheorem{remark}{Remark}
\newcommand{\R}{\mathbb{R}}
\newcommand{\E}{\mathbb{E}}
\DeclareMathOperator{\op}{op}

\title{Synthetic data as vicinal gradient coverage:\\
why smoothness improves while margin falls}
\author{Mechanism theory for the diffusion-data robustness finding (author-derived; verified)}
\date{\today}
\begin{document}
\maketitle

\begin{abstract}
The pilot observation is that adding diffusion-generated data to adversarial training increases the
robust radius even though the raw logit margin \emph{decreases}, because the local sensitivity
$L=\E\|\nabla_x M\|$ falls by more than the margin. This note gives a static mechanism for that
pattern. First-order adversarial training is input-gradient regularization (Lemma~\ref{lem:first-order-at});
synthetic data changes the measure on which the penalty is enforced, adding a positive-semidefinite
\emph{vicinal gradient-coverage matrix}; increasing that matrix suppresses high-gradient modes
(Theorem~\ref{thm:synthetic-filter}); and if a high-gradient mode also carries raw margin, the margin
falls while the sensitivity falls faster, so the scale-invariant radius proxy $M/L$ rises exactly when a
smoother low-gradient component remains (Theorem~\ref{thm:two-mode-smoothness}). A predicted signature,
$L_1/L_2\approx\text{const}$ while both fall (Corollary~\ref{cor:gradient-shape}), is already confirmed in
the pilot ($25.4\to23.6$). The slogan: \emph{diffusion data suppresses high-gradient margin-carrying
modes; the raw margin falls, but the sensitivity falls faster, so $M/L$ rises.}
\end{abstract}

The robust radius admits the first-order decomposition $\Delta\log r\approx\Delta\log M-\Delta\log L$, so a
radius gain with a smaller margin is possible exactly when $L$ falls by more than $M$. The mechanism is not
that synthetic data raises the margin: it is that adversarial training is, to first order, input-gradient
regularization, and synthetic data changes the measure on which this regularizer acts, adding a
positive-semidefinite vicinal gradient-coverage matrix. Modes that carry margin but have high
input-gradient energy are suppressed; the raw margin can fall while the sensitivity falls faster, and the
scale-invariant proxy $M/L$ rises.

\section{First-order adversarial training is input-gradient regularization}

\begin{lemma}[First-order robust-loss expansion]\label{lem:first-order-at}
Let $\ell:\R^d\to\R$ be twice continuously differentiable on the closed ball
$B_\varepsilon(x)=\{x+\delta:\|\delta\|\le\varepsilon\}$ for a norm $\|\cdot\|$ with dual $\|\cdot\|_*$, and
suppose $\|\nabla^2_x\ell(z)\|_{\op}\le H$ for all $z\in B_\varepsilon(x)$ (operator norm induced by
$\|\cdot\|$). Then
\[
\sup_{\|\delta\|\le\varepsilon}\ell(x+\delta)=\ell(x)+\varepsilon\|\nabla\ell(x)\|_*+R_\varepsilon(x),
\qquad |R_\varepsilon(x)|\le\tfrac{H}{2}\varepsilon^2 .
\]
\end{lemma}
\begin{proof}
Set $g=\nabla\ell(x)$. For $\|\delta\|\le\varepsilon$, Taylor with remainder gives
$\ell(x+\delta)=\ell(x)+\langle g,\delta\rangle+\tfrac12\delta^\top\nabla^2\ell(x+s\delta)\delta$ for some
$s\in[0,1]$, and $|\tfrac12\delta^\top\nabla^2\ell(x+s\delta)\delta|\le\tfrac{H}{2}\|\delta\|^2\le\tfrac{H}{2}\varepsilon^2$.
With $\langle g,\delta\rangle\le\|g\|_*\|\delta\|\le\varepsilon\|g\|_*$, the supremum obeys
$\sup_{\|\delta\|\le\varepsilon}\ell(x+\delta)\le\ell(x)+\varepsilon\|g\|_*+\tfrac{H}{2}\varepsilon^2$.
For the lower bound pick $v$ with $\|v\|\le1$ and $\langle g,v\rangle=\|g\|_*$ (the unit ball is compact, so
the linear functional attains its max), and set $\delta=\varepsilon v$. Taylor gives
$\ell(x+\varepsilon v)\ge\ell(x)+\varepsilon\|g\|_*-\tfrac{H}{2}\varepsilon^2$, and the supremum is at least this.
Combining the two bounds gives $|\sup_\delta\ell(x+\delta)-\ell(x)-\varepsilon\|g\|_*|\le\tfrac{H}{2}\varepsilon^2$.
\end{proof}

\begin{remark}[From loss-gradient to margin-gradient]
If $\ell(x,y)=\phi(M(x,y))$ then $\nabla_x\ell=\phi'(M)\nabla_xM$, and on any margin band where
$0<c_-\le|\phi'(M)|\le c_+<\infty$ the loss-gradient and margin-gradient penalties are equivalent up to
constants: $c_-\|\nabla_xM\|\le\|\nabla_x\ell\|\le c_+\|\nabla_xM\|$. Thus the first-order adversarial
penalty controls the same local sensitivity $L=\E\|\nabla_xM\|$ that enters the radius proxy.
\end{remark}

\section{The vicinal gradient-coverage matrix}
Linearize the margin in parameters, $M_\theta(x)=\theta^\top\psi(x)$ with a differentiable feature map
$\psi:\mathcal X\to\R^p$, and write $J_\psi(x)=\nabla_x\psi(x)\in\R^{p\times d}$, so that
$\nabla_xM_\theta(x)=J_\psi(x)^\top\theta$ and $\|\nabla_xM_\theta(x)\|_2^2=\theta^\top J_\psi(x)J_\psi(x)^\top\theta$.
For a (data or vicinal) measure $\nu$ define the gradient-coverage matrix
\[
Q_\nu:=\int J_\psi(x)J_\psi(x)^\top\,d\nu(x),\qquad
\int\|\nabla_xM_\theta(x)\|_2^2\,d\nu(x)=\theta^\top Q_\nu\theta .
\]
$Q_\nu$ is positive semidefinite, since $v^\top Q_\nu v=\int\|J_\psi(x)^\top v\|_2^2\,d\nu(x)\ge0$. Adding
synthetic data changes $\nu$, adding a positive-semidefinite $Q_{\mathrm{syn}}$ to the input-gradient penalty.

\section{Spectral rough-mode filtering}
\begin{theorem}[Synthetic gradient coverage filters high-sensitivity modes]\label{thm:synthetic-filter}
Let $A\in\R^{p\times p}$ be symmetric positive definite, $Q\in\R^{p\times p}$ symmetric positive
semidefinite, $a\in\R^p$, and for $t\ge0$ set $B_t=A+tQ$ and
$\mathcal J_t(\theta)=\tfrac12\theta^\top B_t\theta-a^\top\theta$ with minimizer $\theta_t$. With
$\eta(t)=a^\top\theta_t$ and $S(t)=\theta_t^\top Q\theta_t$,
\[
\eta'(t)=-S(t)\le0,\qquad S'(t)=-2\,\theta_t^\top Q B_t^{-1}Q\theta_t\le0 .
\]
\end{theorem}
\begin{proof}
$B_t$ is symmetric positive definite for $t\ge0$, so $\mathcal J_t$ is strictly convex with unique minimizer
$\theta_t=B_t^{-1}a$ ($B_t\theta_t=a$). Differentiating $B_t\theta_t=a$ in $t$ and using $B_t'=Q$ gives
$Q\theta_t+B_t\theta_t'=0$, i.e.\ $\theta_t'=-B_t^{-1}Q\theta_t$. Then
$\eta'(t)=a^\top\theta_t'=(B_t\theta_t)^\top\theta_t'=\theta_t^\top B_t\theta_t'=\theta_t^\top(-Q\theta_t)=-S(t)$,
and $S(t)=\theta_t^\top Q\theta_t\ge0$. Also
$S'(t)=2(\theta_t')^\top Q\theta_t=2(-B_t^{-1}Q\theta_t)^\top Q\theta_t=-2\theta_t^\top QB_t^{-1}Q\theta_t=-2z^\top B_t^{-1}z$
with $z=Q\theta_t$, and $z^\top B_t^{-1}z\ge0$ since $B_t^{-1}\succ0$.
\end{proof}

\section{A two-mode theorem: margin down, sensitivity down more, radius up}
\begin{theorem}[Smoothness-not-margin in a two-mode model]\label{thm:two-mode-smoothness}
Let $\theta=(\theta_s,\theta_r)\in\R^2$ (smooth mode $s$, rough mode $r$), $A=I_2$,
$Q=\begin{psmallmatrix}0&0\\0&q\end{psmallmatrix}$ with $q>0$, and $a=(a_s,a_r)$ with $a_r\neq0$. For $t\ge0$
let $\theta_t$ minimize $\tfrac12\theta^\top(I_2+tQ)\theta-a^\top\theta$, and set $\eta(t)=a^\top\theta_t$,
$L(t)=\sqrt{\theta_t^\top Q\theta_t}$, $\rho(t)=\eta(t)/L(t)$. Then
\[
\theta_t=\Big(a_s,\ \tfrac{a_r}{1+tq}\Big),\quad
\eta(t)=a_s^2+\tfrac{a_r^2}{1+tq},\quad
L(t)=\tfrac{\sqrt q\,|a_r|}{1+tq},\quad
\rho(t)=\tfrac{(1+tq)a_s^2+a_r^2}{\sqrt q\,|a_r|},
\]
\[
\eta'(t)=-\tfrac{q a_r^2}{(1+tq)^2}<0,\qquad L'(t)=-\tfrac{q\sqrt q\,|a_r|}{(1+tq)^2}<0,\qquad
\rho'(t)=\tfrac{\sqrt q\,a_s^2}{|a_r|}>0\ \ (a_s\neq0),
\]
and, whenever $a_s\neq0$, $-\tfrac{d}{dt}\log L(t)>-\tfrac{d}{dt}\log\eta(t)$. The raw margin decreases, the
sensitivity decreases faster, and the radius proxy $\eta/L$ strictly increases.
\end{theorem}
\begin{proof}
$(I_2+tQ)=\operatorname{diag}(1,1+tq)$, so $(I_2+tQ)\theta_t=a$ gives $\theta_{s,t}=a_s$ and
$\theta_{r,t}=a_r/(1+tq)$. Then $\eta(t)=a_s\cdot a_s+a_r\cdot\tfrac{a_r}{1+tq}=a_s^2+\tfrac{a_r^2}{1+tq}$.
$Q\theta_t=(0,\,q a_r/(1+tq))^\top$, so $\theta_t^\top Q\theta_t=\tfrac{a_r}{1+tq}\cdot\tfrac{q a_r}{1+tq}=\tfrac{q a_r^2}{(1+tq)^2}$
and $L(t)=\sqrt q|a_r|/(1+tq)$ (as $1+tq>0$). Dividing,
$\rho(t)=\big(a_s^2+\tfrac{a_r^2}{1+tq}\big)\cdot\tfrac{1+tq}{\sqrt q|a_r|}=\tfrac{(1+tq)a_s^2+a_r^2}{\sqrt q|a_r|}$.
Differentiating: $\eta'(t)=a_r^2\cdot(-q)(1+tq)^{-2}=-\tfrac{q a_r^2}{(1+tq)^2}<0$;
$L'(t)=\sqrt q|a_r|\cdot(-q)(1+tq)^{-2}=-\tfrac{q\sqrt q|a_r|}{(1+tq)^2}<0$; the numerator of $\rho$ has $t$-derivative
$q a_s^2$ over the constant denominator, so $\rho'(t)=\tfrac{q a_s^2}{\sqrt q|a_r|}=\tfrac{\sqrt q a_s^2}{|a_r|}>0$ when $a_s\neq0$.
For the logarithmic comparison, $\log L(t)=\log(\sqrt q|a_r|)-\log(1+tq)$ gives $-\tfrac{d}{dt}\log L=\tfrac{q}{1+tq}$,
and $-\tfrac{d}{dt}\log\eta=\tfrac{q a_r^2}{(1+tq)^2\eta(t)}$. The inequality $\tfrac{q}{1+tq}>\tfrac{q a_r^2}{(1+tq)^2\eta(t)}$
is, after multiplying by $(1+tq)^2\eta(t)/q>0$, equivalent to $(1+tq)\eta(t)>a_r^2$, i.e.\
$(1+tq)a_s^2+a_r^2>a_r^2$, i.e.\ $(1+tq)a_s^2>0$, i.e.\ $a_s\neq0$.
\end{proof}

\begin{corollary}[Pure temperature shrinkage is the degenerate case]\label{cor:pure-temperature}
If $a_s=0$ in Theorem~\ref{thm:two-mode-smoothness}, then $\rho(t)=|a_r|/\sqrt q$ is constant: suppressing only a
single rough margin-carrying mode scales margin and sensitivity together and cannot improve the radius proxy.
A genuine radius gain requires a remaining smooth component $a_s\neq0$.
\end{corollary}
\begin{proof}
Set $a_s=0$: $\rho(t)=\tfrac{a_r^2}{\sqrt q|a_r|}=\tfrac{|a_r|}{\sqrt q}$, independent of $t$, so $\rho'(t)=0$.
\end{proof}

\begin{corollary}[Uniform gradient-shape shrinkage]\label{cor:gradient-shape}
If the gradient field is dominated by the rough mode, $\nabla_xM_{\theta_t}(x)\approx\theta_{r,t}\,g_r(x)$ with
$g_r$ independent of $t$, then for any two homogeneous norms $\|\cdot\|_\alpha,\|\cdot\|_\beta$,
\[
\frac{\|\nabla_xM_{\theta_t}(x)\|_\alpha}{\|\nabla_xM_{\theta_t}(x)\|_\beta}\approx\frac{\|g_r(x)\|_\alpha}{\|g_r(x)\|_\beta},
\]
approximately independent of $t$: $L_1$ and $L_2$ can both drop sharply while $L_1/L_2$ stays nearly constant.
\end{corollary}
\begin{proof}
Homogeneity gives $\|\nabla_xM_{\theta_t}(x)\|_\alpha\approx|\theta_{r,t}|\,\|g_r(x)\|_\alpha$ and likewise for
$\beta$; since $\theta_{r,t}=a_r/(1+tq)\neq0$, $|\theta_{r,t}|$ cancels in the ratio.
\end{proof}

\section{Interpretation}
The smooth mode $s$ is a low-sensitivity component that generalizes across the vicinity of the data manifold;
the rough mode $r$ is a high-sensitivity component that can raise the raw logit margin on the finite real
training set. Synthetic data increases the vicinal coverage $tq$ on the rough mode, so $\theta_{r,t}=a_r/(1+tq)$
shrinks: the raw margin $\eta(t)=a_s^2+a_r^2/(1+tq)$ falls (the rough mode contributed margin), but the
sensitivity $L(t)=\sqrt q|a_r|/(1+tq)$ falls faster (it is entirely the rough mode), so the radius proxy
$\rho(t)$ rises whenever $a_s\neq0$. A pure logit-temperature rescaling would shrink margin and sensitivity by
the same factor and leave $M/L$ unchanged; that is the degenerate $a_s=0$ case. The observed increase in $M/L$
corresponds to $a_s\neq0$: a smooth component remains while synthetic gradient coverage suppresses rough sensitivity.

\section{Adversarial verification}
\begin{remark}[Scope and falsifiable predictions]\label{rem:av}
\textbf{(1) Not every synthetic data helps.} A useful radius gain needs $Q$ to penalize a rough mode more than a
smooth margin-carrying one; if $Q=qI$ all modes scale uniformly and the mechanism degenerates to
temperature-like shrinkage (no forced gain). \textbf{(2) It explains margin shrink, not denies it:}
$\eta'(t)=-S(t)\le0$. \textbf{(3) It explains $L$ falling faster than $M$:} $-\tfrac{d}{dt}\log L>-\tfrac{d}{dt}\log\eta$
iff $a_s\neq0$, exactly the condition for $\eta/L$ to rise while $\eta$ falls. \textbf{(4) It is not a
robust-overfitting theorem:} the proof uses no time/epochs/learning-rate decay; it is a static
regularization-path result. If the effect appears early and varies monotonically with synthetic diversity, this
is the mechanism; if it appears only after a late robust-overfitting phase, a separate dynamical theorem is
needed. \textbf{(5) Local, not global:} the quadratic model is a local linearization; the theorem identifies the
algebraic condition for the observed decomposition, not a claim that every diffusion-trained net optimizes this
exact objective. \textbf{(6) Falsifiable predictions:} $L$ decreases monotonically with synthetic diversity; $M$
may decrease; $M/L$ increases when a smooth component remains; $L_1/L_2\approx\text{const}$ when the same rough
gradient shape is suppressed in amplitude (Corollary~\ref{cor:gradient-shape}); and the largest $L$ drop should
occur on vicinal/interpolation points where the new synthetic coverage differs most from the empirical measure.
\end{remark}

\paragraph{Empirical status (pilot).} The monotone predictions hold: $L_2$ falls monotonically with synthetic
amount ($5.20\to2.54$), $M$ falls ($4.42\to2.58$), $M/L$ rises, and the gradient-shape prediction
$L_1/L_2\approx\text{const}$ is confirmed ($25.4\to23.6$ while $L_1,L_2$ each roughly halve). The decisive
remaining test is the vicinal-concentration prediction (largest $L$ drop on/near synthetic points), the
\texttt{manifold} measurement.

\end{document}
